O eletrocardiograma reúne informação morfológica, rítmica e espectral em uma
série temporal curta e biologicamente estruturada. Um protocolo automático que
pretende separar batimentos normais de batimentos anômalos precisa preservar o
alinhamento com o R-peak, estabilizar a linha de base e produzir descritores
que continuem legíveis depois da classificação. A MIT-BIH Arrhythmia Database
permanece adequada para esse objetivo por oferecer sinais reais, anotações
públicas e acesso padronizado por WFDB
\citep{moody2001mitdb,goldberger2000physionet,wfdb2022software}.

Morfologia de batimento e intervalos RR seguem entre os descritores mais
interpretáveis em estudos com separação inter-paciente entre treino e teste
\citep{dechazal2004heartbeat,doquire2011featureselection,mousavi2019inter}. O
interesse aqui não é apenas medir uma acurácia global, mas organizar uma
representação compacta que preserve a fisiologia do traçado filtrado e permita
comparar dois regimes de treinamento sob o mesmo vetor de entrada. Por essa
razão, nenhuma feature é extraída do sinal bruto. Toda a modelagem parte de um
ECG previamente filtrado e alinha descritores morfológicos, contexto RR,
informação espectral e respostas Gabor localizadas em regiões aproximadas de P,
QRS e T.

A comparação entre SVM supervisionada e One-Class SVM surge de uma assimetria
recorrente em dados biomédicos. A formulação supervisionada explora exemplos
normais e anômalos no ajuste e funciona como referência para o que o mesmo
conjunto de features consegue oferecer quando há rótulos completos. A
formulação one-class, em contraste, aprende apenas a distribuição dos
batimentos normais e testa se o afastamento dessa região já basta para ranquear
anomalias \citep{cortes1995support,scholkopf2001support,pimentel2014novelty}.
As duas leituras respondem a perguntas diferentes, e a distância entre elas
deve ser interpretada como diferença de informação disponível no treino, não
apenas como troca de algoritmo.

O recorte permanece binário e metodológico. O protocolo não diagnostica classes
clínicas específicas nem tenta reproduzir um fluxo hospitalar completo. O
rótulo \texttt{N} representa normalidade, enquanto os demais símbolos válidos
de batimento entram como anomalia apenas para avaliação. Essa redução concentra
a análise na pergunta de DSP e evita misturar o estudo com classificação
multi-classe de arritmias.

A representação final foi consolidada a partir do notebook do grupo e reúne 18
features filtradas, com seis descritores morfológicos, três features rítmicas,
quatro descritores espectrais e cinco respostas Gabor. A pergunta central passa
a ser direta. Dado esse vetor final, até que ponto a SVM supervisionada e a
One-Class SVM mantêm desempenho útil no split inter-paciente DS1-DS2.

A PR-AUC ocupa o papel central porque a classe anômala é minoritária e o custo
de falsos positivos aparece diretamente na precision \citep{davis2006prroc}.
ROC-AUC, recall, F1 e matrizes de confusão entram como leitura complementar do
regime operacional. O objetivo, portanto, é comparar os dois modelos sob a
mesma representação filtrada final e verificar quanto da separação inter-
paciente depende do acesso prévio a anomalias no treino.
